\section{Discretización y algoritmos}
\label{sec:discretizacion}

\subsection{Balance de energía por celda}
Cada celda activa tiene volumen $V_i=h^3$, caras de área $A_f=h^2$ y distancia $d_i=h/2$ desde el centro a una cara. Al integrar la ecuación de energía en ese volumen y aplicar el teorema de la divergencia, el término conductivo queda expresado como una suma de flujos por caras:
\begin{equation}
  \frac{d}{dt}\int_{V_i}\rho c_pT\,dV
  =\sum_{f\subset\partial V_i}\int_{A_f}
  k\vec{\nabla}T\cdot\vec{n}_{i,f}\,dA
  +\int_{V_i}\dot q_h\,dV-P_{r,i}.
  \label{eq:balance-integral-celda}
\end{equation}
Al aproximar cada flujo por su valor en la cara y la temperatura por el valor del centro, se obtiene
\begin{equation}
  C_i\frac{dT_i}{dt}
  =\sum_{j\in\mathcal N_i}G_{ij}(T_j-T_i)
   +\dot q_{h,i}V_i+Q_{b,i}-P_{r,i},
  \qquad C_i=\rho_i c_{p,i}V_i.
  \label{eq:balance-celda}
\end{equation}
Aquí $\mathcal N_i$ reúne las celdas activas que comparten una cara con $i$. La capacidad $C_i$ se expresa en $\mathrm{J\,K^{-1}}$, la conductancia $G_{ij}$ en $\mathrm{W\,K^{-1}}$, y las potencias de borde entrante $Q_{b,i}$ y de refrigeración $P_{r,i}$ en vatios. Las resistencias de las dos medias celdas se suman en serie:
\begin{equation}
  G_{ij}=\frac{A_{ij}}{d_i/k_i+d_j/k_j}.
  \label{eq:conductancia-vecinos}
\end{equation}
La potencia de borde suma $A_f\phi_f$ sobre las caras exteriores de la celda. En el suelo, la resistencia por unidad de área es $d_i/k_c+L_g/k_g$; en el aire, la resistencia entre el centro y la superficie exterior es $R_{i,f}=d_i/k_c+R_p$. Una cara compartida por dos celdas activas aporta potencias de igual magnitud y signo contrario.

\subsection{Avance temporal y balance radiativo}
El avance explícito de Euler evalúa los intercambios con las temperaturas conocidas al comienzo del paso. Para la hidratación se utiliza directamente el incremento del calor acumulado, que conserva el aporte aun cuando el paso no coincida con el intervalo de las muestras:
\begin{equation}
  T_i^{n+1}=T_i^n+\frac{1}{C_i}\Bigl\{\Delta t_n\Bigl[\sum_{j\in\mathcal N_i}G_{ij}(T_j^n-T_i^n)+Q_{b,i}^n-P_{r,i}^n\Bigr]+V_i\bigl[H_m(a_i^{n+1})-H_m(a_i^n)\bigr]\Bigr\}.
  \label{eq:euler-fundacion}
\end{equation}
En cada cara expuesta se obtiene $T_s$ igualando $(T_s-T_i^n)/R_{i,f}$ con el flujo exterior del balance radiativo. Se emplea Newton con un máximo de nueve iteraciones, una corrección inferior a $10^{-10}\,\mathrm K$ y una comprobación del residuo de $2\times10^{-8}\,\mathrm{W\,m^{-2}}$. La extracción de cada celda se limita a la energía disponible sobre $T_w$ durante el paso antes de aplicar el tope conjunto. El avance del hormigón sigue siendo explícito.

\subsection{Selección del paso y estabilidad}
Para que los coeficientes de una actualización lineal con conductancias positivas sean no negativos, la capacidad de cada celda y su conductancia total de borde $G_{b,i}$ imponen
\begin{equation}
  \Delta t\leq
  \frac{C_i}{\displaystyle\sum_{j\in\mathcal N_i}G_{ij}+G_{b,i}},
  \label{eq:limite-lineal}
\end{equation}
En el problema con radiación, la suma conservadora $S_i$ incluye las conductancias entre vecinos, $A_f/R_{i,f}$ en las caras al aire, las conductancias hacia el suelo y $U_vV_i$ en las celdas refrigerables. El término de aire acota la variación del flujo no lineal respecto de la temperatura interior. El cálculo aplica el siguiente paso:
\begin{equation}
  \Delta t_n=\min\left(
     \Delta t_{\max},\;
     0{,}8\min_{i:S_i>0}\frac{C_i}{S_i},\;
     t_{\mathrm{evento}}-t_n\right).
  \label{eq:paso-utilizado}
\end{equation}
El factor $0{,}8$ deja margen y el último término hace coincidir el avance con activaciones, retiros, cambios horarios y fin del horizonte. La apertura y el cierre se evalúan en los pasos. La restricción de estabilidad se complementó con un contraste temporal para comprobar cuánto cambian los resultados.

\subsection{Temperatura superficial y extracción a \texorpdfstring{$25\,\mathrm{mm}$}{25 mm}}
La primera celda tiene su centro a más de $25\,\mathrm{mm}$ del contorno en ambas mallas. Para estimar la temperatura a esa profundidad se reconstruye un perfil lineal a partir del flujo entrante en cada cara exterior, incluida la base:
\begin{equation}
  T_{b,f}=T_i+\phi_f\frac{d_i}{k_c},
  \qquad
  T_{25,f}=T_i+\phi_f\frac{d_i-\delta}{k_c},
  \qquad \delta=0{,}025\,\mathrm m.
  \label{eq:reconstruccion-s25}
\end{equation}
El desplazamiento $\delta$ sigue la normal interior de la cara cartesiana. Las caras escalonadas explican por qué la posición reconstruida es aproximada en los contornos inclinados. En cada instante se comparan centros y caras para obtener el máximo, y se busca el mínimo entre todos los puntos reconstruidos a $\delta$. El diferencial utiliza ambas temperaturas simultáneas.

\subsection{Resoluciones y sensibilidad de la discretización}
Las dos tasas se compararon con $h=0{,}40\,\mathrm m$ y paso máximo de $300\,\mathrm s$. Para aislar el efecto espacial se refinó la malla de la propuesta manteniendo el origen, las once capas, el paso y la zona refrigerable. En el contraste temporal se mantuvo la malla base y se redujo el paso máximo a la mitad. La Tabla~\ref{tab:contrastes-discretizacion} indica cuántos escenarios cubrió cada prueba.
\begin{table}[htbp]
  \centering
  \caption{Contrastes de sensibilidad utilizados para evaluar la discretización}
  \label{tab:contrastes-discretizacion}
  \small
  \setlength{\tabcolsep}{4pt}
  \begin{tabular}{L{2.0cm}L{3.6cm}L{3.8cm}L{2.5cm}L{2.5cm}}
    \toprule
    Contraste & Configuraciones & Escenarios & Mayor cambio de $T_{\max}$ [K] & Mayor cambio de $\Delta T_{\max}$ [K] \\
    \midrule
    Espacial & $0{,}40\to0{,}20\,\mathrm m$; $300\,\mathrm s$ & $24$ pares de la propuesta & $0{,}669459$ & $0{,}881712$ \\
    Temporal & $300\to150\,\mathrm s$; $0{,}40\,\mathrm m$ & \texttt{M4\_SS} con ERA5-$2024$ y ERA5-$2025$ & $0{,}001148$ & $0{,}001365$ \\
    \bottomrule
  \end{tabular}
  \captionsetup{skip=2pt}
  \caption*{Fuente: elaboración propia a partir de las simulaciones con ambas resoluciones.}
\end{table}

\subsection{Secuencia del algoritmo}
Para seguir una realización basta escoger un plan, una curva y una historia meteorológica. La secuencia se repite con las mismas reglas para todas las combinaciones:
\begin{enumerate}
\item Leer la geometría, preparar la malla y asignar a cada celda su capa y su zona de refrigeración. Cargar el cronograma, las propiedades y los datos del escenario.
\item Comenzar con los circuitos cerrados. Al llegar al inicio de una capa, incorporarla a su temperatura de colocación y actualizar las caras que pasan del aire al contacto con hormigón.
\item Revisar las protecciones según la edad local, calcular el paso admisible y obtener la energía de hidratación generada entre sus dos extremos.
\item Resolver el intercambio de cada cara, comprobar el estado de los circuitos y avanzar simultáneamente las temperaturas. Registrar la energía que entra, sale y queda almacenada.
\item Buscar extremos durante el avance y a ambos lados de los cambios de borde. Guardar sus posiciones e instantes hasta completar el seguimiento.
\item Repetir para los escenarios restantes. Sumar los pesos de los que exceden algún umbral y comparar los planes por su duración y cumplimiento.
\end{enumerate}

\subsection{Equipo y costo computacional}
El equipo identificado al preparar el artículo utilizaba Windows $10$, un Intel Core i7-10750H y $12$ procesadores lógicos. Los tiempos de la Tabla~\ref{tab:tiempos-corridas} proceden de los registros de cada ejecución, también cuando se recuperó un cálculo anterior. La mediana de una simulación fina fue de aproximadamente $12\,\mathrm{min}$.
\begin{table}[htbp]
  \centering
  \caption{Duraciones registradas de las corridas a $672\,\mathrm h$}
  \label{tab:tiempos-corridas}
  \small
  \begin{tabular}{lcccc}
    \toprule
    Configuración & Casos & Mínimo [s] & Mediana [s] & Máximo [s] \\
    \midrule
    $60\,\mathrm{m^3/h}$; $0{,}40\,\mathrm m$; $300\,\mathrm s$ & $24$ & $36{,}8$ & $47{,}8$ & $58{,}3$ \\
    $55\,\mathrm{m^3/h}$; $0{,}40\,\mathrm m$; $300\,\mathrm s$ & $24$ & $55{,}0$ & $85{,}7$ & $96{,}3$ \\
    $60\,\mathrm{m^3/h}$; $0{,}20\,\mathrm m$; $300\,\mathrm s$ & $24$ & $244{,}5$ & $719{,}8$ & $870{,}2$ \\
    $60\,\mathrm{m^3/h}$; $0{,}40\,\mathrm m$; $150\,\mathrm s$ & $2$ & $88{,}9$ & $89{,}9$ & $91{,}0$ \\
    \bottomrule
  \end{tabular}
  \captionsetup{skip=2pt}
  \caption*{Fuente: tiempos conservados en los resultados originales de cada pieza.}
\end{table}
Las $72$ simulaciones finales suman $5{,}243\,\mathrm h$ por proceso, o $5{,}293\,\mathrm h$ al incluir los dos contrastes temporales. Se reutilizaron $26$ resultados y se calcularon $46$ nuevos, que sumaron $4{,}731\,\mathrm h$. La ejecución simultánea impide interpretar esa suma como el tiempo transcurrido entre el comienzo y el final de la campaña.
